\section{Teoria do Modelo de Atores}

O Modelo de Atores define computação como entidades distribuídas (Atores) que se comunicam através de passagem assíncrona de mensagens. Quando um Ator recebe uma mensagem, ele pode concorrentemente:

\begin{enumerate}
    \item Enviar mensagens para endereços de outros Atores que conhece
    \item Criar novos Atores
    \item Designar como lidar com a próxima mensagem que receber
\end{enumerate}

\subsection{Características-Chave}

\begin{itemize}
    \item \textbf{Localidade e Segurança}: Um Ator pode apenas se comunicar com Atores cujos endereços possui, e pode apenas obter endereços através de criação, recepção de mensagem, ou criando novos Atores
    \item \textbf{Comunicação Assíncrona}: Mensagens são enviadas sem aguardar uma resposta, e não há garantia de tempo de entrega
    \item \textbf{Indeterminação}: A ordem de recepção das mensagens não é predeterminada, permitindo modelar concorrência do mundo real
    \item \textbf{Não-determinismo Ilimitado}: Ao contrário de máquinas de Turing, sistemas de Atores podem exibir não-determinismo ilimitado, que melhor representa sistemas distribuídos reais
\end{itemize}

\subsection{Princípios de Sistemas de Informação}

O Modelo de Atores suporta integração robusta de informações através de:

\begin{itemize}
    \item \textbf{Persistência}: Informações são coletadas e indexadas
    \item \textbf{Concorrência}: Trabalho prossegue interativamente e em paralelo
    \item \textbf{Quasi-comutatividade}: Operações podem ser realizadas em ordens diferentes quando a ordem não importa
    \item \textbf{Pluralismo}: Sistemas podem lidar com informações heterogêneas, sobrepostas e inconsistentes
    \item \textbf{Proveniência}: A origem da informação é rastreada e registrada
\end{itemize}
